\documentclass{amsart}

% 数学相关宏包
\usepackage{amsmath,mathtools,amsthm,amssymb}
\usepackage{bbm}
\usepackage{mathrsfs}
\usepackage{centernot}

% 图形和绘图
\usepackage{graphicx,float}
\usepackage{tikz,tikz-cd}
\usepackage{pgfplots}
\usepackage[framemethod=TikZ]{mdframed}
\usetikzlibrary{decorations.pathmorphing,positioning,arrows,automata,calc,shapes,patterns}
\pgfplotsset{compat=1.18}
\usepgfplotslibrary{statistics}

% 表格和列表
\usepackage{booktabs,multirow,colortbl}
\usepackage{enumitem}
\setlist[itemize,enumerate,description]{noitemsep}
\setlistdepth{9}

% 算法
\usepackage{algorithm,algorithmic}

% 字体和颜色
\usepackage{xcolor,listings}
\usepackage{xeCJK}
\setCJKmainfont{Noto Serif CJK SC}

% 页面和链接
\usepackage{fancyhdr,caption}
\usepackage{hyperref}
\usepackage{cleveref}

% 工具宏包
\usepackage{etoolbox,letltxmacro,docmute}
\usepackage{gensymb}

% 自定义设置
\renewcommand{\arraystretch}{1.5}
\let\originalmathbb\mathbb
\renewcommand{\mathbb}[1]{%
    \ifstrequal{#1}{1}{\mathbbm{#1}}{\originalmathbb{#1}}%
}

% 颜色定义
\definecolor{lightblue}{RGB}{173,216,230}
\definecolor{lightgreen}{RGB}{144,238,144}
\definecolor{lightyellow}{RGB}{255,255,224}
\definecolor{lightgray}{RGB}{211,211,211}

		% 超链接设置
		\hypersetup{
		    colorlinks=true,
		    linkcolor=red,
		    filecolor=magenta,
		    urlcolor=cyan,
		    pdftitle={\@title},
		    pdfauthor={\@author},
		    pdfsubject={数学笔记},
		    pdfkeywords={数学},
		    bookmarksnumbered=true,
		    bookmarksopen=true,
		    bookmarksopenlevel=6,
		    pdfstartview=FitH,
		    pdfpagemode=UseOutlines,
		    pdfborder={0 0 0}
		}

% 定理环境
\theoremstyle{plain}
\newtheorem{theorem}{Theorem}[section]
\newtheorem{lemma}[theorem]{Lemma}
\newtheorem{corollary}[theorem]{Corollary}
\newtheorem{proposition}[theorem]{Proposition}

\theoremstyle{definition}
\newtheorem{definition}[theorem]{Definition}
\newtheorem{example}[theorem]{Example}
\newtheorem{exercise}[theorem]{Exercise}
\newtheorem{problem}[theorem]{Problem}

\theoremstyle{remark}
\newtheorem{remark}[theorem]{Remark}
\newtheorem{note}[theorem]{Note}
\newtheorem{observation}[theorem]{Observation}

% 解答环境
\newtheorem*{solution}{Solution}
\newtheorem*{proof*}{Proof}

% 图片路径
\graphicspath{{F:/资料/2025-summer/figures/}}

\author{乐绎华, Yue Yihua}
\address{中山大学, Sun Yat-sen University}
\email{yueyh@mail2.sysu.edu.cn}
\date{\today}
\title{Mathematical Notes}

\begin{document}

\maketitle
\tableofcontents

\chapter{Foundations of Dynamical Systems Theory}

\begin{note}
\textbf{Essential Background for Lyapunov Functions}This document provides the mathematical foundations of dynamical systems theory required to understand Lyapunov functions and their applications to Markov processes.
\end{note}
\section{1. Introduction to Dynamical Systems}

\subsection{1.1 What is a Dynamical System?}

A \textbf{dynamical system} is a mathematical model that describes the time evolution of a system in its state space. It consists of:

\begin{enumerate}
	\item \textbf{State Space} $\mathcal{X}$: The set of all possible states
	\item \textbf{Evolution Rule}: A rule that determines how states change over time
	\item \textbf{Time Domain}: Either discrete ($\mathbb{Z}$ or $\mathbb{N}$) or continuous ($\mathbb{R}$ or $\mathbb{R}_+$)
\end{enumerate}

\subsection{1.2 Types of Dynamical Systems}

\subsubsection{Deterministic Systems}

\begin{itemize}
	\item \textbf{Discrete-time}: $x_{n+1} = f(x_n)$ where $f: \mathcal{X} \to \mathcal{X}$
	\item \textbf{Continuous-time}: $\frac{dx}{dt} = f(x)$ where $f: \mathcal{X} \to \mathbb{R}^d$
\end{itemize}

\subsubsection{Stochastic Systems}

\begin{itemize}
	\item \textbf{Discrete-time}: $X_{n+1} = F(X_n, Z_n)$ where $Z_n$ are random variables
	\item \textbf{Continuous-time}: $dX_t = f(X_t)dt + g(X_t)dW_t$ (stochastic differential equation)
	\begin{itemize}
		\item Here $W_t$ is a \textbf{Wiener process} (Brownian motion), a continuous-time stochastic process with:
		\begin{itemize}
			\item $W_0 = 0$
			\item Independent increments: $W_{t+s} - W_t$ is independent of $\{W_u : u \leq t\}$
			\item Gaussian increments: $W_{t+s} - W_t \sim \mathcal{N}(0, s)$
			\item Continuous paths but nowhere differentiable
		\end{itemize}
	\end{itemize}
\end{itemize}

\subsection{1.3 Key Concepts}

\begin{definition}[Orbit/Trajectory]
The \textbf{orbit} or \textbf{trajectory} of a point $x_0$ is the sequence of states:
	\begin{itemize}
		\item Discrete: $\{x_0, f(x_0), f^2(x_0), \ldots\}$
		\item Continuous: $\{x(t) : t \geq 0\}$ where $x(0) = x_0$
	\end{itemize}
\end{definition}
\begin{definition}[Invariant Set]
A set $S \subseteq \mathcal{X}$ is \textbf{invariant} if trajectories starting in $S$ remain in $S$ for all time.
\end{definition}
\section{2. Stability Theory}

\subsection{2.1 Equilibrium Points}

\begin{definition}[Equilibrium Point]
A point $x^* \in \mathcal{X}$ is an \textbf{equilibrium point} (fixed point) if:
	\begin{itemize}
		\item Discrete: $f(x^*) = x^*$
		\item Continuous: $f(x^*) = 0$
	\end{itemize}
\end{definition}
\subsection{2.2 Types of Stability}

Let $x^*$ be an equilibrium point and consider trajectories starting near $x^*$.

\begin{definition}[Stability (Lyapunov)]
$x^*$ is \textbf{stable} if for every $\varepsilon > 0$, there exists $\delta > 0$ such that:
\[
\|x_0 - x^*\| < \delta \implies \|x_t - x^*\| < \varepsilon \text{ for all } t \geq 0
\]
\end{definition}
\begin{definition}[Asymptotic Stability]
$x^*$ is \textbf{asymptotically stable} if it is stable and additionally:
\[
\lim_{t \to \infty} x_t = x^* \text{ for all } x_0 \text{ in some neighborhood of } x^*
\]
\end{definition}
\begin{definition}[Exponential Stability]
$x^*$ is \textbf{exponentially stable} if there exist constants $C > 0$ and $\alpha > 0$ such that:
\[
\|x_t - x^*\| \leq C\|x_0 - x^*\|e^{-\alpha t}
\]
\end{definition}
\subsection{2.3 Global vs Local Stability}

\begin{itemize}
	\item \textbf{Local stability}: Properties hold in a neighborhood of $x^*$
	\item \textbf{Global stability}: Properties hold for all initial conditions in $\mathcal{X}$
\end{itemize}

\section{3. Lyapunov Functions in Deterministic Systems}

\subsection{3.1 The Lyapunov Function Concept}

\begin{definition}[Lyapunov Function (Deterministic)]
For a dynamical system with equilibrium $x^*$, a function $V: \mathcal{X} \to \mathbb{R}$ is a \textbf{Lyapunov function} if:
	\begin{enumerate}
		\item $V(x^*) = 0$ and $V(x) > 0$ for $x \neq x^*$ (positive definite)
		\item $V$ is continuous
		\item $\frac{dV}{dt} \leq 0$ along trajectories (non-increasing)
	\end{enumerate}
\end{definition}
\subsection{3.2 The Lyapunov Stability Theorem}

\begin{theorem}[Lyapunov's Direct Method]
Consider the system $\dot{x} = f(x)$ with equilibrium $x^* = 0$. If there exists a Lyapunov function $V$ such that:
	\begin{enumerate}
		\item $V(0) = 0$ and $V(x) > 0$ for $x \neq 0$
		\item $\dot{V}(x) = \nabla V(x) \cdot f(x) \leq 0$
	\end{enumerate}
Then the equilibrium is \textbf{stable}.
If additionally $\dot{V}(x) < 0$ for $x \neq 0$, then the equilibrium is \textbf{asymptotically stable}.
\end{theorem}
\textbf{Proof Idea}: The function $V$ acts like an "energy" function that never increases. Level sets $\{x: V(x) = c\}$ form barriers that trajectories cannot cross outward.

\subsection{3.3 Construction of Lyapunov Functions}

\subsubsection{Method 1: Quadratic Forms}

For linear systems $\dot{x} = Ax$, try $V(x) = x^T P x$ where $P > 0$.

The condition $\dot{V} < 0$ becomes:
\[
\dot{V} = x^T(A^T P + PA)x < 0
\]

This requires $A^T P + PA < 0$, which is the \textbf{Lyapunov equation}.

\subsubsection{Method 2: Physical Intuition}

\begin{itemize}
	\item \textbf{Mechanical systems}: Total energy $V = \frac{1}{2}mv^2 + U(x)$
	\item \textbf{Circuit systems}: Total energy stored in capacitors and inductors
	\item \textbf{Economic systems}: Cost or utility functions
\end{itemize}

\subsubsection{Method 3: Gradient Systems}

For systems of the form $\dot{x} = -\nabla U(x)$, the function $V(x) = U(x) - U(x^*)$ is a natural Lyapunov function.

\subsection{3.4 Examples}

\subsubsection{Example 1: Harmonic Oscillator with Damping}

Consider:
\[
\begin{cases}
\dot{x} = y \\
\dot{y} = -kx - cy
\end{cases}
\]
\textbf{Lyapunov function}: $V(x,y) = \frac{1}{2}ky^2 + \frac{1}{2}x^2$ (total energy)

\textbf{Verification}:
\[
\dot{V} = ky\dot{y} + x\dot{x} = ky(-kx - cy) + xy = -cy^2 \leq 0
\]

Since $\dot{V} < 0$ when $y \neq 0$, the origin is asymptotically stable.

\subsubsection{Example 2: Pendulum}

For the damped pendulum $\ddot{\theta} + c\dot{\theta} + \sin\theta = 0$:

\textbf{State variables}: $x_1 = \theta$, $x_2 = \dot{\theta}$
\textbf{System}:
\[
\begin{cases}
\dot{x}_1 = x_2 \\
\dot{x}_2 = -\sin x_1 - cx_2
\end{cases}
\]
\textbf{Lyapunov function}: $V(x_1, x_2) = \frac{1}{2}x_2^2 + (1 - \cos x_1)$

\textbf{Verification}:
\[
\dot{V} = x_2\dot{x}_2 + \sin x_1 \dot{x}_1 = x_2(-\sin x_1 - cx_2) + \sin x_1 x_2 = -cx_2^2 \leq 0
\]

\section{4. Extension to Stochastic Systems}

\subsection{4.1 Stochastic Differential Equations}

Consider the SDE:
\[
dX_t = f(X_t)dt + g(X_t)dW_t
\]

where $W_t$ is a Wiener process.

\subsection{4.2 Itô's Formula}

For a function $V(x,t)$ and the SDE above:
\[
dV(X_t,t) = \left[\frac{\partial V}{\partial t} + \frac{\partial V}{\partial x} \cdot f(X_t) + \frac{1}{2}\text{tr}(g^T(X_t) \nabla^2 V g(X_t))\right]dt + \frac{\partial V}{\partial x} \cdot g(X_t)dW_t
\]

\subsection{4.3 Lyapunov Functions for SDEs}

\begin{definition}[Lyapunov Function (Stochastic)]
A function $V: \mathbb{R}^d \to \mathbb{R}_+$ is a \textbf{Lyapunov function} for the SDE if:
	\begin{enumerate}
		\item $V(x) \geq 0$ with $V(x) = 0$ iff $x = 0$
		\item $\mathcal{L}V(x) \leq 0$ where $\mathcal{L}$ is the generator:
\[
\mathcal{L}V(x) = \frac{\partial V}{\partial x} \cdot f(x) + \frac{1}{2}\text{tr}(g^T(x) \nabla^2 V g(x))
\]
	\end{enumerate}
\end{definition}
\subsection{4.4 Stability in Probability}

\begin{theorem}[Stochastic Stability]
If there exists a Lyapunov function $V$ with $\mathcal{L}V \leq 0$, then the solution is \textbf{stable in probability}.
If additionally $\mathcal{L}V \leq -cV$ for some $c > 0$, then the solution is \textbf{exponentially stable in mean square}.
\end{theorem}
\section{5. Connection to Markov Chains}

\subsection{5.1 Discrete-Time Markov Chains}

A sequence $(X_n)_{n \geq 0}$ is a \textbf{Markov chain} if:
\[
\mathbb{P}(X_{n+1} = j | X_0, X_1, \ldots, X_n) = \mathbb{P}(X_{n+1} = j | X_n)
\]

\subsection{5.2 Markov Operators}

\begin{definition}[Markov Operator]
The \textbf{Markov operator} $P$ acts on functions $f$ as:
\[
(Pf)(x) = \mathbb{E}[f(X_1) | X_0 = x] = \sum_y P(x,y)f(y)
\]
\end{definition}
\subsection{5.3 Lyapunov Functions for Markov Chains}

\begin{definition}[Lyapunov Function (Markov Chain)]
A function $V: \mathcal{X} \to [0,\infty)$ is a \textbf{Lyapunov function} if:
	\begin{enumerate}
		\item $V$ has compact sublevel sets
		\item There exist $\lambda \in (0,1)$ and $b < \infty$ such that:
\[
PV(x) \leq \lambda V(x) + b
\]
	\end{enumerate}
\end{definition}
This is the \textbf{Foster-Lyapunov drift condition}.

\subsection{5.4 Applications to Ergodicity}

\begin{theorem}[Ergodic Theorem]
If a Markov chain has a Lyapunov function, then:
	\begin{enumerate}
		\item The chain has an invariant measure $\pi$
		\item The chain converges to $\pi$ geometrically fast
		\item The invariant measure has finite moments: $\mathbb{E}_\pi[V] < \infty$
	\end{enumerate}
\end{theorem}
\section{6. Mathematical Tools and Techniques}

\subsection{6.1 Supermartingales}

\begin{definition}[Supermartingale]
A sequence $(M_n)$ is a \textbf{supermartingale} if:
\[
\mathbb{E}[M_{n+1} | \mathcal{F}_n] \leq M_n
\]
\end{definition}
\textbf{Connection to Lyapunov functions}: If $V$ is a Lyapunov function for a Markov chain, then $V(X_n)$ is a supermartingale.

\subsection{6.2 Optional Stopping Theorem}

\begin{theorem}[Optional Stopping]
For a supermartingale $(M_n)$ and a stopping time $\tau$:
\[
\mathbb{E}[M_\tau] \leq \mathbb{E}[M_0]
\]
\end{theorem}
\textbf{Application}: Gives bounds on expected hitting times and exit probabilities.

\subsection{6.3 Concentration Inequalities}

\subsubsection{Azuma-Hoeffding Inequality}

For a martingale $(M_n)$ with bounded differences:
\[
\mathbb{P}(M_n - M_0 \geq t) \leq e^{-t^2/(2\sum_{i=1}^n c_i^2)}
\]

\subsubsection{Chernoff Bound}

For sums of independent random variables:
\[
\mathbb{P}(S_n \geq t) \leq e^{-\lambda t}\mathbb{E}[e^{\lambda S_n}]
\]

\subsection{6.4 Functional Analysis Tools}

\subsubsection{Banach Spaces}

\begin{itemize}
	\item \textbf{Completeness}: Every Cauchy sequence converges
	\item \textbf{Norms}: $\|f\|_\infty = \sup_x |f(x)|$, $\|f\|_p = (\int |f|^p)^{1/p}$
\end{itemize}

\subsubsection{Operators}

\begin{itemize}
	\item \textbf{Bounded operators}: $\|Pf\| \leq C\|f\|$
	\item \textbf{Compact operators}: Map bounded sets to relatively compact sets
	\item \textbf{Spectral theory}: Eigenvalues and eigenfunctions
\end{itemize}

\section{7. Advanced Topics}

\subsection{7.1 Generators of Markov Processes}

For a continuous-time Markov process $(X_t)$:
\[
(\mathcal{L}f)(x) = \lim_{t \to 0} \frac{1}{t}[\mathbb{E}^x[f(X_t)] - f(x)]
\]

\subsection{7.2 Resolvent and Potential Theory}

The \textbf{resolvent} operator:
\[
R_\lambda = (\lambda I - \mathcal{L})^{-1}
\]

\textbf{Green's function}: $G(x,y) = \int_0^\infty p_t(x,y)dt$

\subsection{7.3 Large Deviations}

\begin{theorem}[Cramér's Theorem]
For i.i.d. random variables with moment generating function $M(\lambda)$:
\[
\lim_{n \to \infty} \frac{1}{n}\log \mathbb{P}(S_n \geq na) = -\sup_\lambda(\lambda a - \log M(\lambda))
\]
\end{theorem}
\subsection{7.4 Ergodic Theory}

\subsubsection{Birkhoff's Ergodic Theorem}

For an ergodic transformation $T$:
\[
\lim_{n \to \infty} \frac{1}{n}\sum_{k=0}^{n-1} f(T^k x) = \int f d\mu
\]

\subsubsection{Mixing Properties}

\begin{itemize}
	\item \textbf{Strong mixing}: $\lim_{n \to \infty} \mathbb{P}(A \cap T^{-n}B) = \mathbb{P}(A)\mathbb{P}(B)$
	\item \textbf{Weak mixing}: $\lim_{n \to \infty} \frac{1}{n}\sum_{k=0}^{n-1} |\mathbb{P}(A \cap T^{-k}B) - \mathbb{P}(A)\mathbb{P}(B)| = 0$
\end{itemize}

\section{8. Connections to Other Areas}

\subsection{8.1 Optimization Theory}

\textbf{Gradient descent}: $x_{k+1} = x_k - \alpha \nabla f(x_k)$

\textbf{Lyapunov function}: $V(x) = f(x) - f(x^*)$

\subsection{8.2 Control Theory}

\textbf{Optimal control}: Minimize $J = \int_0^\infty L(x,u)dt$ subject to $\dot{x} = f(x,u)$

\textbf{Lyapunov-based control}: Design $u$ so that $V$ is a Lyapunov function for the closed-loop system.

\subsection{8.3 Machine Learning}

\textbf{Stochastic gradient descent}: $\theta_{k+1} = \theta_k - \alpha \nabla L(\theta_k, \xi_k)$

\textbf{Convergence analysis}: Use Lyapunov functions to prove convergence to minima.

\section{9. Computational Aspects}

\subsection{9.1 Numerical Methods}

\subsubsection{Euler-Maruyama Method}

For SDE $dX_t = f(X_t)dt + g(X_t)dW_t$:
\[
X_{n+1} = X_n + f(X_n)\Delta t + g(X_n)\Delta W_n
\]

\subsubsection{Milstein Method}

Higher-order scheme:
\[
X_{n+1} = X_n + f(X_n)\Delta t + g(X_n)\Delta W_n + \frac{1}{2}g(X_n)g'(X_n)[(\Delta W_n)^2 - \Delta t]
\]

\subsection{9.2 Simulation Techniques}

\subsubsection{Monte Carlo Methods}

\begin{itemize}
	\item \textbf{Importance sampling}: Change of measure to reduce variance
	\item \textbf{Markov Chain Monte Carlo}: Sample from complex distributions
	\item \textbf{Multilevel Monte Carlo}: Reduce computational cost
\end{itemize}

\subsubsection{Quasi-Monte Carlo}

Use low-discrepancy sequences instead of random numbers.

\subsection{9.3 Verification of Lyapunov Conditions}

\subsubsection{SOS (Sum of Squares) Programming}

Express $V$ and $-\dot{V}$ as sums of squares of polynomials.

\subsubsection{Semidefinite Programming}

Relaxation of polynomial optimization problems.

\section{10. Historical Notes and Further Reading}

\subsection{10.1 Historical Development}

\begin{itemize}
	\item \textbf{1892}: A.M. Lyapunov introduces the direct method
	\item \textbf{1931}: A.N. Kolmogorov formalizes probability theory
	\item \textbf{1953}: K.L. Chung develops Markov chain theory
	\item \textbf{1960s}: R.L. Dobrushin and others develop ergodic theory
	\item \textbf{1993}: S.P. Meyn and R.L. Tweedie unify the theory
\end{itemize}

\subsection{10.2 Key References}

\subsubsection{Foundational Texts}

\begin{enumerate}
	\item \textbf{Khalil, H.K.} (2002). \textit{Nonlinear Systems} (3rd ed.). Prentice Hall.
	\item \textbf{Meyn, S.P. \& Tweedie, R.L.} (2009). \textit{Markov Chains and Stochastic Stability}. Cambridge University Press.
	\item \textbf{Øksendal, B.} (2003). \textit{Stochastic Differential Equations}. Springer.
\end{enumerate}

\subsubsection{Advanced Topics}

\begin{enumerate}
	\item \textbf{Ethier, S.N. \& Kurtz, T.G.} (1986). \textit{Markov Processes: Characterization and Convergence}. Wiley.
	\item \textbf{Revuz, D. \& Yor, M.} (1999). \textit{Continuous Martingales and Brownian Motion}. Springer.
	\item \textbf{Durrett, R.} (2019). \textit{Probability: Theory and Examples}. Cambridge University Press.
\end{enumerate}

\subsection{10.3 Modern Applications}

\begin{itemize}
	\item \textbf{Financial mathematics}: Risk management, portfolio optimization
	\item \textbf{Biology}: Population dynamics, epidemiology
	\item \textbf{Engineering}: Control systems, signal processing
	\item \textbf{Computer science}: Algorithm analysis, machine learning
	\item \textbf{Physics}: Statistical mechanics, quantum systems
\end{itemize}

\section{Summary}

This foundation provides the essential mathematical background for understanding:

\begin{enumerate}
	\item \textbf{Deterministic stability theory} and classical Lyapunov functions
	\item \textbf{Stochastic processes} and their stability properties
	\item \textbf{Markov chain theory} and the Foster-Lyapunov drift conditions
	\item \textbf{Connections between} different areas of mathematics and applications
\end{enumerate}

The key insight is that Lyapunov functions provide a unified framework for analyzing stability and convergence across deterministic and stochastic systems, with applications ranging from classical mechanics to modern machine learning algorithms.

\textbf{Next Steps}: With this foundation, you can now dive deeper into the specific techniques for finding Lyapunov functions for Markov chains, as covered in your other lecture notes.


\end{document}